% TOP_PROBLEM: {"id": 3102, "title": "Pebbling a cartesian product", "category_id": 3, "category": {"id": 3, "name": "graph_theory", "display_name": "Graph Theory", "description": "Problems involving graphs, networks, and their properties.", "slug": "graph-theory", "order_index": 3, "created_at": "2026-07-31T15:26:25.671Z"}, "status": "open", "problem_number": "OPG-586", "set_id": 12}
% FIRST_POSTED: 2026-10-01
% ATTEMPT_STATUS: unresolved
% COMPLETION: 5
% MODEL: GPT 6 Astra Ultra
% UnsolvedMath ID 3102; source catalogue code OPG-586
% !TeX program = lualatex
% Research attempt, 2026-10-01. Canonical statement preserved verbatim.
\documentclass[11pt,a4paper]{article}
\usepackage[margin=25mm]{geometry}
\usepackage{fontspec}
\setmainfont{lmroman10-regular.otf}[BoldFont=lmroman10-bold.otf,
 ItalicFont=lmroman10-italic.otf,BoldItalicFont=lmroman10-bolditalic.otf]
\usepackage{amsmath,amssymb,mathtools}
\usepackage{unicode-math}
\setmathfont{latinmodern-math.otf}
\AtBeginDocument{\let\setminus\smallsetminus}
\usepackage{xurl,hyperref}
\hypersetup{colorlinks=true,urlcolor=blue}
\setcounter{secnumdepth}{0}
\setlength{\parindent}{0pt}
\setlength{\parskip}{5pt}
\setlength{\emergencystretch}{3em}
\begin{document}
\section*{Pebbling a cartesian product}\label{3102}
{\small UnsolvedMath ID 3102 · OPG-586 · Graph Theory · OpenGarden}
\subsection{Definitions and mathematical statement}
Let $G$ be a finite nonempty connected simple undirected graph.
A pebble configuration is a function $C:V(G)\to\mathbb Z_{\ge0}$,
with total size $\sum_v C(v)$. A legal pebbling move chooses adjacent
vertices $u,v$ with $C(u)\ge2$, removes two pebbles from $u$, and adds
one pebble at $v$. One pebble is therefore lost at each move.

For a chosen target vertex $t$, a configuration is solvable at $t$ if
some finite sequence of legal moves leaves at least one pebble there.
The pebbling number $p(G)$ is the least integer $N$ such that every
configuration of total size $N$ is solvable at every target vertex.
The sequence of moves may depend on both the configuration and the
target; simultaneous delivery to all targets is not required.

For two such graphs, their Cartesian product $G_1\mathbin\Box G_2$ has
vertex set $V(G_1)\times V(G_2)$. Two distinct pairs are adjacent if
one coordinate is equal and the other coordinate is an adjacent pair
in its factor graph.
Graham's conjecture asserts
\[
 p(G_1\mathbin\Box G_2)\le p(G_1)p(G_2)
\]
for every pair of finite connected simple graphs.

Thus the product bound must handle arbitrary, possibly highly uneven,
placements on the product graph. A construction that succeeds only for
placements distributed uniformly among factor fibres is insufficient.
Connectedness ensures that pebbling numbers are finite; the single-vertex
graph is permitted and has pebbling number one.
\subsection{Short English statement}
Is the pebbling number of a Cartesian product at most the product of the pebbling numbers of its factors?
\subsection{Research attempt}
\paragraph{Literature check and attempted route.}
Chung's original pebbling work [A1] and the original problem page [S2] were
checked on 1 October 2026. Products of paths are known positive cases; the
present proof does not claim those results as new. We test the most direct
fibre-by-fibre strategy for general factors, derive a rigorous bound with
an explicit loss, and show on four vertices why rounding away the residues
cannot prove Graham's sharp product bound.

\paragraph{Collecting disjoint packets in a fibre.}
Write $a=p(G)$, $b=p(H)$ and $h=|V(H)|$, and target $(g_0,h_0)$.
For each $y\in V(H)$ let $c_y$ be the number of pebbles in the $G$-fibre
$V(G)\times\{y\}$. Partition any $a\lfloor c_y/a\rfloor$ of those pebbles
into packets of size $a$. By definition of $p(G)$, each packet can supply a
pebble to $(g_0,y)$ using moves inside the fibre. Treat packets as coloured
only for bookkeeping: process one packet at a time, never consuming the
pebbles reserved from another. Every move legal for the selected packet is
legal for the actual uncoloured configuration. Thus at least
$\lfloor c_y/a\rfloor$ pebbles can be collected at $(g_0,y)$.
If their sum over $y$ is at least $b$, the definition of $p(H)$ finishes the
job along the fibre $\{g_0\}\times V(H)$.

\paragraph{A general upper bound with a controlled error.}
Write $c_y=aq_y+r_y$, $0\le r_y\le a-1$. If $\sum_yq_y\le b-1$, then
\[
 N=\sum_yc_y\le a(b-1)+h(a-1).
\]
Consequently every configuration of size
\[
 a(b-1)+h(a-1)+1=ab+(a-1)(h-1)
\]
is solvable. Interchanging the two factors gives the proved estimate
\[
 p(G\mathbin\Box H)\le p(G)p(H)+
 \min\{(p(G)-1)(|V(H)|-1),(p(H)-1)(|V(G)|-1)\}.
\]
The single-vertex case makes the error zero. In general the error is genuine
in this argument, even though it is not evidence that the conjecture is false.
It represents packets that cannot be completed using only their initial fibre.

\paragraph{A four-vertex stress test of the sharp inference.}
For $G=H=K_2$, each factor has pebbling number two: a target already occupied
needs no move, and otherwise both pebbles at its neighbour give one move.
Use coordinates $(i,j)\in\{0,1\}^2$, target $(0,0)$, and initial counts
\[
 C(0,0)=0,\ C(1,0)=1,\ C(0,1)=1,\ C(1,1)=2.
\]
There are four pebbles, exactly the conjectural product. The two fibre totals
are one and three, and their rounded packet counts sum to only
$\lfloor1/2\rfloor+\lfloor3/2\rfloor=1<2$.
Nevertheless the configuration is solvable: move the two pebbles from
$(1,1)$ to $(0,1)$, which now has two, then move them to $(0,0)$.
The successful route uses a residue already sitting on the collecting fibre.
Discarding residues loses precisely that interaction. The example does not
refute fibre methods; it refutes the claim that the packet lower bound alone
always supplies $b$ pebbles at size $ab$.

\paragraph{Verification and remaining gap.}
The packet proof checks legality, order of processing, and arbitrary uneven
initial distributions. The algebraic threshold follows from an integer
contrapositive, so the final $+1$ is necessary. A finite exact search over
all 35 distributions of four pebbles on the four-cycle verifies solvability
at a fixed target; symmetry covers all targets. It also verifies the displayed
two-move certificate directly. For general factors a sharp proof must exploit
residues while preserving the ability to gather enough pebbles; the current
bound deliberately leaves their possible contribution unproved.

\paragraph{Reproducibility.}
The finite calculations above are implemented in
\texttt{scripts/check\_attempt\_batch\_2026\_10\_01.py}; run it with Python 3
from the repository. It uses exact integer arithmetic and no external packages.

\paragraph{Final status.}
\textbf{UNRESOLVED.} A general bound with an additive error and an exact
obstruction to a naive sharp proof are established. The product inequality
itself is not proved for arbitrary factors. Completion estimate: 5\%.

\subsection{Sources}
\begin{itemize}
\item[S1] ULAM.AI and the UnsolvedMath Contributors, supplied problems.json, record 3102 (OPG-586), OpenGarden collection. Catalogue identity and source transcription; CC BY 4.0.
\url{https://huggingface.co/datasets/ulamai/UnsolvedMath}.
\item[S2] Open Problem Garden, Pebbling a cartesian product. Original problem and discussion; the mathematical formulation below is expanded from the supplied source record.
\url{http://www.openproblemgarden.org/op/pebbling_a_cartesian_product}.
\item[A1] F. R. K. Chung, \emph{Pebbling in Hypercubes}, SIAM Journal on Discrete Mathematics 2 (1989), 467--472, primary publisher abstract and bibliographic record. \url{https://doi.org/10.1137/0402041}.
\end{itemize}
\end{document}
